\documentclass[11pt]{article}
\usepackage[letterpaper,left=1in,right=1in,top=.88in,bottom=.88in]{geometry}
\usepackage{amsmath,amssymb,amsthm,booktabs,microtype}
\usepackage[hidelinks]{hyperref}
\newtheorem{theorem}{Theorem}[section]
\newtheorem{lemma}[theorem]{Lemma}
\newtheorem{corollary}[theorem]{Corollary}
\newcommand{\ii}{\mathrm i}
\newcommand{\dd}{\mathrm d}
\newcommand{\Var}{\operatorname{Var}}
\newcommand{\E}{\mathbb E}
\newcommand{\ind}{\mathcal I_N}
\hypersetup{pdftitle={Strict total positivity of the Hessian of optimal zero-cancelling filters}}
\title{Strict total positivity of the Hessian of optimal zero-cancelling filters}
\author{}\date{}
\begin{document}
\maketitle
\begin{abstract}
A positive average that removes prescribed Fourier frequencies has a least
possible support width. We determine its full quadratic response when a
multiple frequency splits. In the odd coefficients of the centered root
polynomial, every minor of the width Hessian is strictly positive.
We give a finite formula for every entry and an explicit orthogonal
decomposition. It determines the leading cost along every asymmetric
splitting direction, including directions where the first six or ten
orders are invisible. The proof identifies the perturbations of the
optimal atoms with Bernoulli cumulants, diagonalizes the variance of their
displacements in shifted Legendre coordinates and proves strictness by
counting disjoint paths. On the critical line this gives the local cost
of cancelling clusters of Dirichlet zero terms with positive prime averages.
\end{abstract}

\section{The extremal problem and the result}
For a compactly supported probability measure $\mu$ on $\mathbb R$ let
$F_\mu(\lambda)=\int e^{\ii\lambda t}\dd\mu(t)$. Given a multiset of $N$
positive real frequencies, let $B$ be the smallest support diameter among
measures whose transform vanishes at every prescribed frequency, to the
order of its multiplicity. Translation preserves the constraints.
Normalize the mean $\kappa=N^{-1}\sum_j\lambda_j$ and set
\[
 u_j=\lambda_j/\kappa-1,\qquad
 P(x)=\prod_j(x-u_j)=x^N+\sum_{r=2}^Nc_rx^{N-r},
 \qquad b=\kappa B .
\]
Thus $c_r=(-1)^r e_r(u)$, where $e_r$ is the elementary symmetric
function, and the missing coefficient $c_1$ is zero. Write $E$ for the
even-indexed coefficients, $O$ for the odd-indexed ones, and
$\ind=\{3,5,\ldots\}\cap[3,N]$. Assume $N\ge3$.

At coalescence an optimal measure in normalized time has atoms at
$j\pi$ with weights $2^{-N}\binom Nj$, $0\le j\le N$. The ProveIt
article establishes a unique nearby optimizing branch and the estimate
$b=N\pi+O(\max|u_j|^3)$ \cite{repo}. A subsequent research packet
proves $b-N\pi\asymp_N\|O\|^2$, evaluates the coefficient on the
$c_3$ axis for all $N$, and gives one $c_5$ example \cite{previous}.
It leaves open the evaluation of the complete quadratic form. We
show that this form has a strict sign constraint on every minor.

\begin{theorem}[Width Hessian]\label{thm:main}
There is a symmetric matrix $H_N$, indexed by $\ind$, such that
near coalescence
\[
 b(E,O)=N\pi+O^{\mathsf T}H_NO+
 O_N\bigl((\|E\|+\|O\|)\|O\|^2\bigr).
\]
The notation $O^{\mathsf T}H_NO$ denotes the quadratic form in the
vector $O$. Twice $H_N$ is the odd-coordinate Hessian. For all
increasing index sets $I,J\subseteq\ind$ of equal positive size,
\[
                     \det (H_N)_{I,J}>0 .
\]
Define the Bernoulli cumulants $K_r(p)$ by
$\log(1-p+pe^t)=\sum_{r\ge1}K_r(p)t^r/r!$. Put
\[
 M_r(p)=\frac{(-1)^{(r-1)/2}}{(r-1)!}\, K_r(p),
 \qquad M'_r(p)=\sum_{n\in\ind}a_{nr}P_{n-1}(2p-1),
\]
where $P_j$ is the ordinary Legendre polynomial normalized by $P_j(1)=1$.
For $n>r$ we have $a_{nr}=0$. Then
\begin{align}
 (H_N)_{rt}&=\pi^{r+t-1}\sum_{n\in\ind}\beta_{N,n}a_{nr}a_{nt},\label{eq:factor}\\
 \beta_{N,n}&=
 \frac{(N+n-1)!(N-n)!/((N-1)!N!)-1}{(2n-1)n(n-1)}.\label{eq:beta}
\end{align}
All constants are explicit. With $\left\{\begin{smallmatrix}r\\ k\end{smallmatrix}\right\}$
a Stirling number of the second kind, for $n\le r$,
\[
 a_{nr}=\frac{(2n-1)(-1)^{(r-1)/2}}{(r-1)!}
 \sum_{k=n}^r (-1)^{k-1}k!\left\{\begin{matrix}r\\ k\end{matrix}\right\}
 \frac{((k-1)!)^2}{(k-n)!(k+n-1)!}.
\]
\end{theorem}

In particular each pair of asymmetric coefficient perturbations has a
positive mixed quadratic cost when their signs agree.
Every inverse entry has strict sign $(-1)^{i+j}$ when $i,j$ denote
positions in the increasing list $\ind$, by the cofactor formula.
Both phenomena concern the interaction between perturbations and go
beyond positive definiteness. These signs persist for the actual Hessian
$\partial^2_{OO}b(E,O)$ on a neighborhood of the origin, since there
are finitely many strictly positive minors and the analytic continuation
constructed below has continuous derivatives. In particular, for two
distinct odd coordinate directions $i,j$ and positive increments $h,k$,
\[
 \begin{split}
 &b(E,O+h\mathbf e_i+k\mathbf e_j)-b(E,O+h\mathbf e_i)\\
 &\hspace{30mm}{}-b(E,O+k\mathbf e_j)+b(E,O)>0
 \end{split}
\]
whenever the rectangle lies in a sufficiently small neighborhood. This
follows by integrating the positive mixed derivative over the rectangle.
If its four corners have real-rooted frequency polynomials, each term
is a genuine optimal width; the analytic continuation suffices on the
interior.

\begin{corollary}[Every asymmetric ray]\label{cor:ray}
For a fixed centered real vector $v$ whose multiset lacks reflection
symmetry about zero, let $r$ be its least odd index with $e_r(v)\ne0$.
For $u_j=hv_j$ with $h\to0$,
\[
 b-N\pi=(H_N)_{rr}e_r(v)^2h^{2r}+O_{N,v}(h^{2r+2}).
\]
\end{corollary}
Indeed $E=O_{N,v}(h^2)$, the first odd coefficient is $-e_r(v)h^r$,
and the next possible one has degree $r+2$. A real centered multiset
has all odd coefficients zero precisely when its roots are invariant
under reflection.

\section{An integral for the quadratic form}
We give the analytic reduction to specify exactly where optimality and
the quadratic approximation enter. This reduction also appears in
\cite{previous}. Normalize $\kappa=1$ and anchor the first
atom at zero. By the splitting theorem of \cite{repo}, the global
optimizer at real nearby frequencies has $N+1$ atoms. Its width is its
last node. At $P=x^N$ the weights are $w_j^0=2^{-N}\binom Nj$.
The nodes $t_j$ and weights $w_j$ extend real analytically in the
coefficient variables $c$ to a neighborhood of the origin. In detail,
for a holomorphic $f$, the polynomial remainder of degree less than $N$
modulo $P$ is given by the contour integral
\[
 R_Pf(x)=\frac1{2\pi\ii}\oint
 \frac{f(z)}{P(z)}\frac{P(z)-P(x)}{z-x}\,\dd z
\]
on a fixed circle enclosing all roots. Set the real and imaginary parts
of the coefficients of $R_P(F_\mu(1+\cdot))$ to zero and impose unit
mass. At the base the linearized transform has the form
$A(z)+(\ii+\ell/\pi)T(z)$, with $z=e^{\ii\pi x}$,
$\ell=\log z$ near $z=1$, and $A,T$ real polynomials of degree at
most $N$. Here $T(0)=0$ from anchoring and $A(-1)=0$ from unit mass.
Vanishing through degree $N-1$ in $z-1$ first forces $T$ divisible
by $(z-1)^N$ by imaginary parts, then $A$ by real parts. Their extra
conditions force $A=T=0$. All weights at the base are positive, so the
square real Jacobian is invertible. The implicit function theorem
provides the continuation, and uniqueness identifies it with the
source optimizer on admissible real-rooted data.

Let $G(s)=\sum_j w_j e^{s t_j}$. For all sufficiently small
coefficient vectors its only zeros in $0<\Im s<2$ are
$\ii(1+u_j)$, with orders. At the base $G(s)=((1+e^{\pi s})/2)^N$.
Its two endpoint terms dominate uniformly at the two horizontal ends of
the strip, also after perturbation, since the node gaps and endpoint
weights stay bounded below. On the remaining compact rectangle Rouch\'e's
theorem preserves the zero count. Divisibility by $P$ identifies all
$N$ zeros, also on the continuation where some $u_j$ may be complex.
Their imaginary parts sum to $N$.

For completeness, integration of
$\Delta\log |G|=2\pi\sum_{G(\rho)=0}\delta_\rho$
in horizontal strips gives, with $b=t_N$,
\[
 J(T):=\int_{\mathbb R}\log\frac{G(x)}{|G(x+\ii T)|}\dd x
 =\frac{bT^2}{2}-2\pi\sum_{0<\Im\rho<T}(T-\Im\rho).
\]
Indeed the integrated horizontal second derivative of $\log|G|$
is the difference $b$ of the endpoint slopes. Thus
$J''=b-2\pi\sum\delta_{\Im\rho}$ for $T>0$, with
$J(0)=J'(0)=0$. Endpoint domination gives exponential tails for the
logarithmic difference, and local zero singularities are integrable.
At $T=2$ this becomes
\begin{equation}\label{eq:strip}
 b-N\pi=\frac12\int_{\mathbb R}
 -\log\left|\sum_j p_j(x)e^{2\ii t_j}\right|\dd x,
 \qquad p_j(x)=w_je^{xt_j}/G(x).
\end{equation}
It is nonnegative throughout the real coefficient neighborhood.
When $O=0$, nearby real-rooted mirror configurations give width
$N\pi$ by the exact reflection theorem of \cite{repo}. Distinct pairs
of real roots (and a zero root for odd $N$) give an open set in the
even coefficients by the simple-root theorem. Analyticity extends
$b(E,0)=N\pi$ throughout the even coefficient neighborhood. Thus
$\partial_O b(E,0)=0$ by the local minimum in \eqref{eq:strip}.

For an odd coefficient variation $v$, denote its node velocities
at the origin by $s_j$. Expansion of the integrand in \eqref{eq:strip}
gives the quadratic coefficient
\begin{equation}\label{eq:variance}
 Q(v)=\frac1\pi\int_0^1\frac{\Var(s_K)}{p(1-p)}\dd p,
 \qquad K\sim\operatorname{Bin}(N,p).
\end{equation}
Indeed the coefficient of the squared perturbation in the negative
log-modulus is twice the variance of the $s_j$ under the base weights
$p_j(x)$. Derivatives of the weights sum to zero and do not enter
this coefficient. Factoring out the phase of the dominant endpoint
bounds the differentiated integrands on each tail by
$C(1+|x|)^2 e^{-g|x|}$ for a fixed gap $g>0$. At the base the weights
are binomial with $p=e^{\pi x}/(1+e^{\pi x})$, giving the stated
Jacobian. Taylor integration in $O$ at fixed $E$ and analyticity
give the remainder of Theorem~\ref{thm:main}.

\section{The velocities are Bernoulli cumulants}
For a unit variation in $c_r$, $r$ odd, at the origin define
$T_r(z)=\sum_j(-1)^jw_j^0s_j z^j$ and $C=(-1)^N2^{-N}$.
At base the transform evaluated at $1+x$ is $C(z-1)^N$.
Differentiate its divisibility by $x^N+c_r x^{N-r}$. Modulo $x^N$
the first variation is $C(z-1)^N/x^r$. Substituting
$x=\ell/(\ii\pi)$ and taking imaginary parts yields
\begin{equation}\label{eq:rem}
 T_r(z)\equiv (-1)^{(r-1)/2}\pi^r C\frac{(z-1)^N}{(\log z)^r}
               \pmod{(z-1)^N},\qquad T_r(0)=0 .
\end{equation}
This determines a polynomial of degree at most $N$ uniquely.

Define the Eulerian numerator $D_{r-1}(z)$ by
$\sum_{m\ge1}m^{r-1}z^m=D_{r-1}(z)/(1-z)^r$ for $|z|<1$.
It vanishes at zero and has degree at most $r-1$ since $r\ge3$.
To see its jet at $1$, apply $(\dd/\dd y)^{r-1}$ to
$e^y/(1-e^y)=-1/y+\text{holomorphic}$ at $y=0$.
For odd $r$ this gives
\[
 \frac{D_{r-1}(z)}{(r-1)!}
 =\frac{(z-1)^r}{(\log z)^r}+O((z-1)^r).
\]
It follows that $T_r$ equals
$(-1)^{(r-1)/2}\pi^r C(z-1)^{N-r}D_{r-1}(z)/(r-1)!$.
The binomial mean of the velocities is
$2^N(1-p)^N T_r(-p/(1-p))$.
Writing $z=-p/(1-p)$, the defining generating function for cumulants
gives $K_r(p)=-D_{r-1}(z)/(1-z)^r$. This follows first for $p<1/2$
by expanding $\log(1-p+pe^t)$ in powers of $z e^t$, then for all
$p$ by polynomial identity. Consequently
\begin{equation}\label{eq:means}
            \E[s_K]=\pi^r M_r(p).
\end{equation}
This step gives a single family of polynomials for every cluster size
$N\ge r$. The values at $p=0,1$ vanish and $M_r(1-p)=-M_r(p)$.
Hence $M'_r$ is symmetric about $1/2$ with zero integral. Its
Legendre indices are among the even degrees $2,4,\ldots,r-1$.

\section{Orthogonal decomposition of the variance}
If $f$ is any real function on $\{0,\dots,N\}$ and
$m(p)=\E f(K)$, then, writing $(N)_s=N!/(N-s)!$,
\begin{equation}\label{eq:hoeffding}
 \Var f(K)=\sum_{s=1}^N\frac{[p(1-p)]^s}{s!(N)_s}
                   (m^{(s)}(p))^2 .
\end{equation}
Write $K=\sum_i X_i$ for independent Bernoulli variables of mean $p$.
The multilinear extension of $f(\sum_i X_i)$ on the cube, expanded
at $(p,\dots,p)$, has common coefficient $m^{(s)}(p)/(N)_s$ for
each $\prod_{i\in S}(X_i-p)$ with $|S|=s$. These products are orthogonal,
which proves the identity.

If $m'(p)=\sum_{n}h_nP_{n-1}(2p-1)$, the $s$th term of
\eqref{eq:hoeffding} after integration in \eqref{eq:variance}
is diagonal in the $h_n$. In fact for $1\le s\le n$,
\[
 \int_0^1[p(1-p)]^{s-1}
 \left(\frac{\dd^{s-1}}{\dd p^{s-1}}P_{n-1}(2p-1)\right)^2\dd p
 =\frac{(n+s-2)!}{(2n-1)(n-s)!}.
\]
Mixed degrees are orthogonal. These identities follow by
$s-1$ integrations by parts in the Legendre equation, or from the
associated Legendre norm after $x=2p-1$. Hence the coefficient of
$h_n^2/\pi$ is
\[
 \frac1{2n-1}\sum_{s=1}^{n}\frac{(n+s-2)!}{s!(n-s)!(N)_s}.
\]
Include $s=0$ and multiply the sum inside the prefactor by
$n(n-1)$. The summands become
$\binom ns\binom{n+s-2}s/\binom Ns$. The terminating
Chu--Vandermonde identity evaluates this sum as
$(-N-n+1)^{\overline n}/(-N)^{\overline n}$,
the factorial ratio in \eqref{eq:beta}.
Here bars denote rising factorials and $n\le N$ ensures all terms
through the terminating index are defined.
Subtracting the unit $s=0$ term proves \eqref{eq:beta}.
Equations \eqref{eq:means}--\eqref{eq:hoeffding} yield \eqref{eq:factor}
by polarization. To obtain $a_{nr}$, expand
\[
 K_r(p)=\sum_{k=1}^r (-1)^{k-1}(k-1)!
                \left\{\begin{matrix}r\\ k\end{matrix}\right\}p^k
\]
in the logarithmic generating function. Rodrigues' formula gives
$\int_0^1p^{k-1}P_{n-1}(2p-1)\dd p=((k-1)!)^2/((k-n)!(k+n-1)!)$
for $k\ge n$, and zero at smaller $k$. This proves the claimed sum.
Jacobi diagonalizations for related Bernstein--Durrmeyer operators
are classical \cite{waldron}. Here the cumulant map identifies which
perturbations of the extremizing atoms enter those modes.

\section{Strict positivity of every minor}
The cumulants satisfy $K_1(p)=p$ and
$K_{r+1}(p)=p(1-p)K'_r(p)$. For example this follows by
differentiating $\log(1+e^\theta)$ in the natural parameter
$\theta=\log(p/(1-p))$. On $f=K'_r$, the resulting transition is
$f\mapsto (p(1-p)f)'$. With $x=2p-1$ this operator sends $P_l$ to
\begin{equation}\label{eq:walk}
 \frac{l(l-1)}{2(2l+1)}P_{l-1}
 -
 \frac{(l+1)(l+2)}{2(2l+1)}P_{l+1}.
\end{equation}
For $l=0$ omit the first term. Formula \eqref{eq:walk} follows from the
Legendre derivative identity and the three-term recurrence for $xP_l$.

Use a directed lattice with time increasing by one and height $l\ge0$.
Give up and down steps the absolute values in \eqref{eq:walk};
delete zero-weight steps. The positive down steps start at height at
least two. The coefficient of $P_{n-1}$ in $K'_r$ is the signed sum
of the paths from $(1-r,0)$ to $(0,n-1)$. Every path has exactly
$(r+n-2)/2$ up steps. The sign of $a_{nr}$ for $n\le r$ therefore
equals $(-1)^{(n-1)/2}$ after applying the sign in $M_r$.
More precisely, $(r-1)!\,|a_{nr}|$ is the unsigned path weight.

Take increasing lists $r_1<\cdots<r_k$ and $n_1<\cdots<n_k$ from
$\ind$. The determinant of the unsigned path matrix, with sources in
the order $r_i$ and sinks in the order $n_j$, is a sum of weights
of vertex-disjoint families connecting corresponding indices.
Indeed, cancellation of intersecting families follows by swapping the
tails at a first intersection in the determinant expansion
\cite{gv}. A path starting earlier from height zero has strictly
positive height when a later path starts, since no downward step to
height zero has positive weight. If it ends below the later path, the
two must meet. Their heights at common times have the same parity,
so they cannot cross in the middle of edges. Therefore only the
order-preserving matching contributes.

If $n_i\le r_i$ for every $i$, a disjoint family exists. For each $i$
take all the up steps first, followed by all its down steps. Its height
at time $\tau$ in $[1-r_i,0]$ is
$\min(\tau+r_i-1,n_i-1-\tau)$. At times shared by two such paths,
both arguments of the minimum for the higher-index path exceed those
of the lower-index one by at least two. The paths are disjoint and
all steps have positive weight. If some $n_i>r_i$, the
order-preserving connection is impossible by path length. We have proved
\[
 \det(|a_{n_i r_j}|)_{i,j=1}^k\ge0,\qquad
 \det(|a_{n_i r_j}|)_{i,j=1}^k>0
 \ \Longleftrightarrow\ n_i\le r_i\text{ for all }i .
\]
The factors $(r-1)!$ are positive, and transposition leaves each
determinant unchanged.

Absorb the powers of $\pi$ in \eqref{eq:factor} into positive
diagonal scalings. Since the signs of the coefficients only depend
on $n$, the remaining matrix is
$|a|^{\mathsf T}\operatorname{diag}(\beta_{N,n})|a|$.
Every $\beta_{N,n}$ is strictly positive. Cauchy--Binet expands
each minor in nonnegative terms. With a minor of size $k$, take
$n_i=2i+1$ for $1\le i\le k$. Both minors in this term are strictly
positive, since the $i$th index of either given list is at least
$2i+1$. This proves the strict inequality in Theorem~\ref{thm:main}.

\section{An exceptional cluster and the role of GRH}
For $N=7$ let $D=\operatorname{diag}(\pi^3,\pi^5,\pi^7)$ in the order
$3,5,7$. The exact normalized matrix is
\[
 \pi D^{-1}H_7 D^{-1}=
 \begin{pmatrix}
 7/600 & 1/1440 & 7/216000\\
 1/1440 & 11/17280 & 371/5702400\\
 7/216000 & 371/5702400 & 189091/1010880000
 \end{pmatrix}.
\]
For instance the minor on rows $3,5$ and columns $5,7$ compares two
mixed responses to a third, a relation invisible if only one
coefficient varies at a time.

Let $v_1,\ldots,v_7$ be the roots of
$x(x^2-1)(x^2-4)(x^2-9)+1$. They are real and simple. At either end
of each interval $[j-1/4,j+1/4]$, $j=-3,\ldots,3$, the absolute
value of the unshifted product exceeds one. Indeed it is at least
$(1/4)(3/4)^2(7/4)^2(11/4)^2>1$, since each positive integer
distance from $j$ occurs at most twice among the other roots.
The unshifted product changes sign across each interval, and adding one
preserves these seven sign changes. The centered polynomial has
$c_3=c_5=0,\ c_7=1$. Consequently the seven frequencies
$\lambda_j=\kappa(1+hv_j)$ satisfy
\[
 \kappa B-7\pi=
 \frac{189091\,\pi^{13}}{1010880000}\,h^{14}+O(h^{16}).
\]
This gives a family whose leading width response occurs at order fourteen.

A zero term $x^{\beta+\ii\gamma}$ in a prime explicit formula acquires
the multiplier $\int e^{(\beta+\ii\gamma)t}\dd\mu(t)$ when the
prime statistic is averaged at scales $xe^t$. Tilting by
$e^{\beta t}$, followed by normalization, identifies its zeros and
their orders with the Fourier conditions at $\gamma$. The tilt is
invertible and preserves support. Thus for prescribed Dirichlet zero
ordinates on the critical line the theorem evaluates the width Hessian
of a positive cancelling average. The exact formulas apply to arbitrary
real clusters. Whether actual zero ordinates realize the exceptional
algebraic example is a separate arithmetic question. The filter
theorem is unconditional. The arithmetic application uses the common
real part supplied by the repository's Dirichlet GRH input.

\paragraph{Reproducibility and scope.}
The accompanying Python script uses exact rational arithmetic to compare
the factorization against direct covariance integration after computing
node velocities by interpolation. It checks every diagonal and mixed
entry for $3\le N\le9$ and verifies the signs of all their minors. The arguments above cover every $N$.
The result is local at a coalesced frequency. Computing the Hessian at
general mirror configurations with a positive optimal measure, and
determining where its minors can change sign, remain open directions.
Optimal input shaping provides a separate engineering context for the
underlying minimum-duration problem \cite{lau}.

\begin{thebibliography}{9}
\bibitem{repo} ProveIt contributors, \emph{Beyond the Riemann hypothesis for
Dirichlet $L$-functions}, section \emph{Positive filters that cancel zeros},
especially the reflection and splitting theorems.
\href{https://github.com/VladimirReshetnikov/ProveIt/blob/4f99495218d1a19087deabd8fb5bbb8f5df72a46/Analysis/GeneralizedRiemannHypothesis/Article/sections/fil.tex}{Pinned source, commit 4f99495218}.
\bibitem{previous} \emph{The sixth-order cost of splitting a zero-cancelling
filter}, companion research packet, November 2026,
\texttt{grh\_cayley\_and\_filters.zip}, file
\texttt{filters/cluster\_filters.tex}. The final paragraph asks for the
uniform mixed entries of the quadratic form.
\bibitem{waldron} S. Waldron, \emph{A generalised beta integral and the limit
of the Bernstein--Durrmeyer operator with Jacobi weights}, J. Approx.
Theory 122 (2003), 141--150.
\href{https://doi.org/10.1016/S0021-9045(03)00041-8}{Publisher record}.
\bibitem{gv} I. Gessel and G. Viennot, \emph{Binomial determinants, paths,
and hook length formulae}, Adv. Math. 58 (1985), 300--321.
\bibitem{lau} M. A. Lau and L. Y. Pao, \emph{Input shaping and time-optimal
control of flexible structures}, Automatica 39 (2003), 893--900.
\href{https://doi.org/10.1016/S0005-1098(03)00024-4}{Publisher record}.
\end{thebibliography}
\end{document}
